\documentclass{article}
\usepackage{graphicx} % Required for inserting images
\usepackage{hyperref} % Required for hyper links
\usepackage[a4paper, margin=2.5cm]{geometry} % For easier layout setting
\usepackage{fancyhdr} % For fancy header
\usepackage{lastpage} % To find last page
\usepackage{amsmath}

% The instruction is written in A4 vertical page so that
% instructor can show this on the side while showing what to do in the browser at the class.
% and students can do that as well since they have to work on their laptop without extra monitor.

\pagestyle{fancy}
\graphicspath{ {./images/} }
\graphicspath{ {../images/} }

\title{2026 DMSC Summer School Exercise Student Guide}
\author{ESS DMSC Staff}
\date{24 August 2026}

\fancyhead[L]{2026 DMSC Summer School Exercise Student Guide}
\fancyhead[C]{}
\fancyhead[R]{24 August 2026}

% Use header and skip making title.
\fancyfoot[C]{Page \thepage\ of \pageref{LastPage}}
\def\_{{\tt \char`\_ }} % underscore is not defined as a plain character.

\begin{document}
\sffamily  % Change font
\begin{center} % Insert ESS logo
\makebox{
  \raisebox{-\totalheight}[0pt][0pt]{
    \includegraphics[width=6cm]{ess-logo-opacity-10}
  }
}
\end{center}

Before we start, make sure you have joined our discord channel.

The link should be in the welcome email.

\section{Sign in at ESS VISA/ Create an instance}

\subsection{ESS Account}

\begin{description}

    \item[If you don’t have one yet:] see the 'How to make an ESS account' guide we shared in the welcome email.
    \item[If you just created the ESS account:] it will not work up to an hour, so wait until your ESS account is properly set up in our system or go to the \nameref{sec:local-env} section instead.
    \item[ORCID ID] is required to sign in at VISA. E-mail does not work for logging in at VISA.

\end{description}

\subsection{Sign in at VISA}
\begin{description}
    \item[Go to \href{https://visa.ess.eu}{https://visa.ess.eu}] and sign in with the ORCID ID.
\end{description}

\subsection{Create a new instance}
\begin{description}

    \item[Search for experiments]
    \item[Select a proposal] under 'DMSC Summer School 2026' - Proposal: 919936

    \includegraphics[width=7cm]{visa-1.png}

    \item[Check] the 'Terms and conditions'
    \item[Create] the instance and wait until jupyter lab button shows up. It may take up to a few minutes.

    \includegraphics[width=8cm]{visa-1-2.png}

\end{description}

\section{Download Workbook}

\subsection{Open Jupyter Lab of the VISA Instance}
\includegraphics{visa-2.png}

\subsection{Open a terminal}
\includegraphics{visa-3.png}

\subsection{Download summer school workbook}

The following command will clone the repo and make a symlink to your home folder.

\begin{description}

    \item[Command:] source /ess/data/workshop/919936/derived/tools/initialize\_visa.sh
    \item[Link to the workbook:] \href{https://github.com/ess-dmsc-dram/dmsc-school-notebooks}{github.com/ess-dmsc-dram/dmsc-school-notebooks}
    \item[Feel free to fork] the repo and add your own remote repo to the local git configuration! *Note that git command only works in a notebook as a magic command or in a terminal.
    \item[After these steps,] you should be able to see the workbook folder in your home like the following pictures. \includegraphics{visa-4.png} \includegraphics{visa-5.png}

\end{description}

\section{Warning!}
\subsection{Non-continuous Autosaving of Notebooks}
Unfortunately, the notebook status will not be stored in your browser if you escape the VISA jupyter page.
Changing tabs are okay, but closing tab or going to home in the VISA page will remove the progress.
However, kernels will still be alive, keeping the memory alive. Make sure you save the notebook regularly to avoid such confusing state. Auto-save is enabled but the interval is 120s by default.

\subsection{Regular Log-out from VISA}
VISA platform logs users out regularly for security reasons, but it may not refresh the page automatically.
Therefore, if the jupyter lab is not responding expectedly, try refreshing the VISA page or re-login.

\appendix
\section{Useful Links}

Useful links to documentation or web services for the school exercises.

\subsection{\href{https://ess-dmsc-dram.github.io/dmsc-school/intro.html}{DMSC Summer School Textbook}}

https://ess-dmsc-dram.github.io/dmsc-school/intro.html : There are hints in the textbook webpage. Bookmark it and feel free to see the hints anytime you feel stuck.

\subsection{\href{https://visa.ess.eu/}{VISA}}

https://visa.ess.eu/ : Virtual computing platform for ESS.

\subsection{\href{https://indico.ess.eu/event/3961/overview}{Indico page}}

https://indico.ess.eu/event/3961/overview: Indico page to the summer school.

\subsection{\href{https://staging.scicat.ess.eu/}{Scicat Web UI (staging)}}

https://staging.scicat.ess.eu/ : Dataset catalouging system.

\subsection{\href{https://scipp.github.io/}{Scipp Documentation}}

https://scipp.github.io/: Data reduction python API documentation page.


\section{Create local environment}

\label{sec:local-env}
If you don't have the ESS account yet, you may need to create a local environment.

\subsection{Using Pixi}
\begin{description}

    \item[Install Pixi:] https://pixi.prefix.dev/latest/installation/
    \item[Clone the Workbook:] git clone https://github.com/ess-dmsc-dram/dmsc-school-notebooks.git
    \item[Change Working Directory:] cd dmsc-school-notebooks
    \item[Install Virtual Environment:] pixi install \text{-}\text{-}frozen
    \item[Start Jupyter Lab Locally:] pixi run \text{-}\text{-}frozen jupyter lab .

\end{description}

\subsection{Using Conda}

\begin{description}

    \item[Clone the Workbook:] git clone https://github.com/ess-dmsc-dram/dmsc-school-notebooks.git
    \item[Download Lock File:] https://git.esss.dk/dram/code-shelf/code-shelf/-/packages/343
    \item[Create Conda Environment:] conda env create -n dmsc-ss \text{-}\text{-}file path-to-the-downloaded-lock-file.yml
    \item[Activate Conda Environment:] conda activate dmsc-ss
    \item[Start Jupyter Lab Locally:] jupyter lab path-to-the-cloned-workbook

\end{description}

\begin{center} % Insert ESS logo
\makebox{
  \raisebox{0pt}[0pt][0pt]{
    \includegraphics[width=6cm]{ess-logo-opacity-10}
  }
}
\end{center}
\end{document}
